\documentclass[11pt]{article}
\usepackage{ap-psych}
\title{Unit 1: Biological Bases of Behavior: Question Set}
\author{Paul Wang}
\date{November 2026}
\docpart{Unit 1 \textperiodcentered\ Question Set}

\begin{document}
\maketitle

\begin{abstract}
Exam-style multiple choice on all six Unit 1 topics, followed by a full Evidence-Based Question on sleep and adolescent learning. Most items present a scenario or study and ask you to name the structure, chemical, process, or conclusion it shows; two items are built on a shared data display. All items and sources are original and written for practice; the EBQ sources are fictional summaries modeled on the exam's format. Allow about 30 minutes for the multiple choice and 45 minutes for the EBQ.
\end{abstract}

\scoringline{MCQ: \blank\ / \qtotal \qquad EBQ: \blank\ / 7 \qquad Target: 21+ MCQ and 5+ EBQ \qquad Pace: 72 s per MCQ}

\section{Heredity and the nervous system (1.1--1.2)}

\begin{q}{2}{1.1 Evolutionary psych}{A preference that aided ancestral survival is an evolutionary explanation.}
Which explanation of humans' strong preference for sweet and fatty foods reflects evolutionary psychology?
\ch{Children learn to like sweets because parents use them as rewards.}{Ancestors who sought calorie-dense foods were more likely to survive and reproduce.}{Advertising creates cravings for sugar.}{Sweet foods trigger the release of leptin.}
\end{q}

\begin{q}{4}{1.1 Adoption studies}{Resemblance to biological (not adoptive) parents implicates genes.}
Adopted children's scores on a measure of extraversion correlate \(r = .35\) with their biological parents' scores and \(r = .04\) with their adoptive parents'. This suggests extraversion is
\ch{mostly determined by shared home environment}{unrelated to genes or environment}{caused by the adoption process}{influenced by genetic factors}
\end{q}

\begin{q}{1}{1.1 Epigenetics}{Environment changing gene expression without altering DNA = epigenetics.}
Rat pups who receive more licking and grooming from their mothers show changes in which stress-related genes are active, without any change to their DNA sequence. This illustrates
\chh{epigenetics}{eugenics}{heritability}{natural selection}
\end{q}

\begin{q}{3}{1.2 Autonomic NS}{After the threat passes, the parasympathetic system restores calm.}
Ten minutes after a near-miss car accident, Dana's heart rate slows and her breathing returns to normal. Which division is now dominant?
\chh{Sympathetic}{Somatic}{Parasympathetic}{Central}
\end{q}

\begin{q}{2}{1.2/1.3 Reflex arc}{Reflexes route through spinal interneurons; the brain is not required.}
A patient with a spinal cord injury that blocks signals to the brain still jerks his leg when his knee is tapped. This is possible because
\ch{the brain processes the tap through the thalamus}{a sensory neuron, interneuron, and motor neuron in the spinal cord complete the response}{the parasympathetic nervous system controls the leg}{the cerebellum is intact}
\end{q}

\section{Neurons, neurotransmitters, hormones, drugs (1.3)}

\begin{q}{1}{1.3 Neural firing}{Sodium influx makes the inside more positive: depolarization.}
When a neuron reaches threshold, positively charged ions rush into the axon. This process is called
\chh{depolarization}{reuptake}{the refractory period}{the resting potential}
\end{q}

\begin{q}{4}{1.3 Reuptake}{Blocking reuptake leaves serotonin in the synapse longer.}
SSRIs, a common antidepressant class, increase serotonin activity by
\ch{destroying serotonin receptors}{increasing the speed of action potentials}{acting as serotonin antagonists}{blocking the reabsorption of serotonin by the sending neuron}
\end{q}

\begin{q}{3}{1.3 Neurotransmitters}{Tremors + difficulty starting movement = Parkinson's, low dopamine.}
A 62-year-old develops a hand tremor and has trouble initiating movement. A drug that increases which neurotransmitter would most likely be prescribed?
\chh{Serotonin}{GABA}{Dopamine}{Substance P}
\end{q}

\begin{q}{2}{1.3 Neurotransmitters}{GABA is inhibitory; boosting it calms neural activity and reduces anxiety.}
Antianxiety drugs that enhance the effects of the brain's main inhibitory neurotransmitter target
\chh{glutamate}{GABA}{norepinephrine}{acetylcholine}
\end{q}

\begin{q}{1}{1.3 Opioids}{Opioids mimic endorphins at their receptors (agonists).}
Morphine relieves pain because it
\ch{acts as an agonist at endorphin receptor sites}{blocks dopamine reuptake}{stimulates the release of substance P}{increases melatonin production}
\end{q}

\begin{q}{4}{1.3 Withdrawal}{Withdrawal symptoms are often opposite to the drug's effect: a stimulant's crash brings fatigue and low mood.}
After quitting a stimulant, a person feels exhausted and depressed for several days. These symptoms reflect
\chh{tolerance}{an agonist effect}{synesthesia}{withdrawal}
\end{q}

\begin{q}{3}{1.3 Hormones}{Oxytocin promotes bonding and trust.}
Which hormone is most closely associated with social bonding between a parent and newborn?
\chh{Leptin}{Adrenaline}{Oxytocin}{Ghrelin}
\end{q}

\section{The brain (1.4)}

\begin{q}{2}{1.4 Brain}{Balance and coordination = cerebellum (also alcohol-sensitive).}
A field sobriety test asks drivers to walk heel-to-toe and touch their nose. It primarily checks the functioning of the
\chh{hypothalamus}{cerebellum}{amygdala}{occipital lobe}
\end{q}

\begin{q}{4}{1.4 Cortex}{Planning and impulse control = prefrontal cortex (cf.\ Phineas Gage).}
After a workplace accident damages his prefrontal cortex, a careful accountant becomes impulsive and makes reckless decisions. This mirrors the case of
\ch{H.M.\ and the hippocampus}{a split-brain patient}{Broca's aphasia}{Phineas Gage}
\end{q}

\begin{q}{1}{1.4 Cortex}{Body-touch map in the parietal lobe = somatosensory cortex.}
Stimulating a point on the strip of cortex at the front of the parietal lobe makes a patient feel a tingling in her thumb. The researcher stimulated the
\ch{somatosensory cortex}{motor cortex}{visual cortex}{auditory cortex}
\end{q}

\begin{q}{3}{1.4 Plasticity}{Reassignment of function after damage, strongest in childhood = neuroplasticity.}
A 4-year-old has much of her left hemisphere removed to stop seizures; within a few years her language skills are nearly normal. This illustrates
\chh{lateralization}{lesioning}{neuroplasticity}{the all-or-nothing principle}
\end{q}

\begin{q}{2}{1.4 Limbic system}{Hunger, thirst, temperature regulation = hypothalamus.}
Damage to which structure would most likely disrupt a rat's regulation of body temperature and eating?
\chh{Thalamus}{Hypothalamus}{Corpus callosum}{Medulla}
\end{q}

\section{Sleep (1.5)}

\begin{stimulus}
\textbf{Figure.} A sleep lab records one participant's night. Minutes spent in each stage per 90-minute cycle:
\begin{center}
\begin{tabular}{lccccc}
\toprule
 & Cycle 1 & Cycle 2 & Cycle 3 & Cycle 4 & Cycle 5\\\midrule
NREM-3 & 40 & 30 & 10 & 0 & 0\\
REM & 10 & 15 & 25 & 30 & 40\\
\bottomrule
\end{tabular}
\end{center}
\end{stimulus}

\begin{q}{4}{1.5 Sleep cycles}{The table shows REM rising and deep NREM-3 falling across the night.}
Which conclusion is supported by the figure?
\ch{Deep sleep increases as the night goes on.}{REM sleep is constant across cycles.}{Sleep cycles shorten to 45 minutes by morning.}{REM periods lengthen and deep sleep disappears later in the night.}
\end{q}

\begin{q}{1}{1.5 Sleep cycles}{Vivid dreams are most likely where REM is longest (Cycle 5).}
Based on the figure, the participant is most likely to report a vivid dream if woken during
\chh{Cycle 5}{Cycle 1}{Cycle 2}{Cycle 3}
\end{q}

\begin{q}{3}{1.5 REM rebound}{After REM deprivation, the body compensates with extra REM.}
After two nights in which researchers woke her whenever REM began, a participant sleeps undisturbed and shows far more REM than usual. This is
\chh{insomnia}{narcolepsy}{REM rebound}{sleep apnea}
\end{q}

\begin{q}{2}{1.5 Circadian rhythm}{Night-shift work misaligns the light-set biological clock.}
A nurse who switches to night shifts struggles to sleep during the day and feels groggy at work. The best explanation is a disruption of her
\ch{REM rebound}{circadian rhythm}{vestibular sense}{reticular activating system only}
\end{q}

\section{Sensation (1.6)}

\begin{q}{1}{1.6 Vision}{Rods dominate the periphery and work in dim light.}
Astronomers look slightly to the side of a faint star to see it better. This works because
\ch{rods, concentrated in the periphery, are more sensitive in dim light}{cones in the fovea detect dim light best}{the blind spot is in the center of the retina}{the lens cannot accommodate at night}
\end{q}

\begin{q}{4}{1.6 Hearing}{High pitches are coded by the place of maximum vibration on the basilar membrane.}
Which theory best explains how we perceive a very high-pitched whistle?
\ch{Frequency theory}{Opponent-process theory}{Gate-control theory}{Place theory}
\end{q}

\begin{q}{3}{1.6 Adaptation}{Constant stimulus \(\rightarrow\) receptors reduce response = sensory adaptation.}
After a few minutes in a friend's house, Ana no longer notices the smell of their dog. This is
\chh{transduction}{sensory interaction}{sensory adaptation}{Weber's law}
\end{q}

\begin{q}{2}{1.6 Smell}{Smell bypasses the thalamus with direct links to the amygdala and hippocampus.}
The smell of sunscreen instantly brings back an emotional memory of a childhood beach trip. This is best explained by the fact that smell
\ch{is processed first in the occipital lobe}{has direct pathways to emotion and memory areas such as the amygdala and hippocampus}{relies on rods and cones}{is relayed through the thalamus like all other senses}
\end{q}

\begin{q}{1}{1.6 Body senses}{Sensing body-part position and movement = kinesthesis.}
With her eyes closed, a dancer can tell exactly where her arms and legs are. This relies on
\chh{kinesthesis}{the vestibular sense}{gustation}{blindsight}
\end{q}

\section{Evidence-Based Question}

\begin{frq}{Question 2 \textperiodcentered\ EBQ \textperiodcentered\ 7 points \textperiodcentered\ 45 minutes}
\small
\textbf{Directions.} Read the three summarized sources (fictional, written for practice) and answer the prompt. Cite sources by author and year.

\textbf{Prompt.} Using evidence from at least two sources, develop and justify an argument about whether later school start times improve adolescents' academic functioning.

\medskip
\textbf{Source 1.} Rivera \& Tan (2023). \emph{Participants}: 2{,}400 students at 8 U.S.\ high schools. \emph{Method}: Correlational. Researchers compared sleep logs and grades at four schools that moved start times from 7:30 to 8:45 a.m.\ with four schools that did not. \emph{Results}: Students at later-start schools slept an average of 42 more minutes per night; sleep duration correlated \(r = .31\) with GPA. First-period tardiness fell by 30\%.

\textbf{Source 2.} Lindqvist (2022). \emph{Participants}: 64 Swedish adolescents aged 15--17. \emph{Method}: Experiment. Participants were randomly assigned to 9 or 5 hours in bed for five nights in a sleep lab, then completed a word-pair learning task. \emph{Results}: The 9-hour group recalled 24\% more word pairs (\(p < .01\)). EEG showed more slow-wave (NREM-3) sleep in the 9-hour group.

\textbf{Source 3.} Mensah, Ortiz, \& Park (2024). \emph{Participants}: 410 adolescents in Ghana and Mexico. \emph{Method}: Survey. Students reported bedtime, wake time, and daytime sleepiness. \emph{Results}: Melatonin release (estimated from self-reported sleepiness onset) occurred about 1--2 hours later than in adults. Students with school before 7:30 a.m.\ reported the highest daytime sleepiness.

\medskip
\begin{enumerate}[label=(\Alph*),leftmargin=2em,itemsep=3pt]
\item Propose a specific and defensible claim based in psychological science that responds to the prompt.
\item (i) Support your claim using at least one piece of specific and relevant evidence from one of the sources. (ii) Explain how the evidence from Part B(i) supports your claim using a psychological perspective, theory, concept, or research finding learned in AP Psychology.
\item (i) Support your claim using an additional piece of specific and relevant evidence from a different source. (ii) Explain how the evidence from Part C(i) supports your claim using a different psychological perspective, theory, concept, or research finding learned in AP Psychology.
\end{enumerate}
\end{frq}

\makekey

\subsection*{EBQ scoring guide (practice rubric, 7 points)}
\small
Modeled on the published EBQ structure (claim, two pieces of cited evidence, reasoning with course content).
\begin{compare}{Part}{Earns credit}{Common miss}
A. Claim (1) & ``Starting high school later improves adolescents' learning because it lets them get more sleep that matches their delayed circadian rhythm.'' & ``Sleep is important''; restating the prompt as a question. \\
B(i). Evidence (1) & Lindqvist (2022): randomly assigned 9-hour sleepers recalled 24\% more word pairs than 5-hour sleepers. & Vague: ``a study showed sleep helps.'' \\
B(ii). Reasoning (2) & 1 pt: explains how the evidence supports the claim (more sleep \(\rightarrow\) better learning). 1 pt: uses a course concept accurately, e.g., \emph{memory consolidation theory}: extra NREM-3 sleep strengthens newly learned material; experiment allows a causal claim. & Defines consolidation without linking it to the 24\% result. \\
C(i). Evidence (1) & Mensah et al.\ (2024): adolescents' melatonin release occurs 1--2 hours later than adults', and the earliest start times produced the most sleepiness. \emph{Or} Rivera \& Tan (2023): later start = 42 more minutes of sleep; sleep correlated \(r = .31\) with GPA. & Using a second fact from the same source as B. \\
C(ii). Reasoning (2) & 1 pt: links evidence to claim. 1 pt: different concept, e.g., the \emph{circadian rhythm} (suprachiasmatic nucleus and melatonin shift later in adolescence, so early starts force waking during biological night); or correlation limits (\(r = .31\) is moderate and cannot prove causation, but converges with Lindqvist's experiment). & Claiming Rivera \& Tan \emph{proves} later starts cause higher GPA. \\
\end{compare}

\end{document}
